% !TeX encoding = UTF-8
\documentclass[12pt, a4paper]{article}

\usepackage{geometry}
\usepackage{amsmath, amsfonts, amssymb, amsthm} % стандартный набор AMS-пакетов для математ. текстов
\usepackage{mathtext}
\usepackage[utf8]{inputenc} % кодировка utf8
\usepackage[russian]{babel} % русский язык
\usepackage[pdftex,dvipsnames]{xcolor} % работа с цветами
\usepackage[pdftex]{graphicx} % графика (картинки)
\usepackage{tikz,pgfplots} % рисунки
\usepackage{booktabs}
\usepackage{indentfirst}
%\usepackage[labelfont=bf,labelsep=endash,skip=3pt]{caption} % подпись картинок
% \usepackage{fancyhdr,pageslts} % настройка колонтитулов
\usepackage{enumitem} % работа со списками
\usepackage{floatrow,multicol,multirow,longtable,hhline} % работа с таблицами
\usepackage{float,wrapfig} % плавающие объекты
\usepackage{tcolorbox} % рамка вокруг текста
%\usepackage[calc]{datetime2} % дата
\usepackage{bm} % жирное начертание в формулах
\usepackage{physics} % физический пакет
\DeclareMathAlphabet\mathbfcal{OMS}{cmsy}{b}{n}
\usepackage{pgfornament} % красивые рюшечки и вензеля
\usepackage{mdframed}
\usepackage{derivative}
\usepackage{mathrsfs} %EDS
\usepackage{soul} % strikethorugh

%\usepackage{boondox-cal}

% ----------------------------------------
% Настройка шрифта

% Просто закооментируйте следующую строчку, если не работает. Будет другой шрифт, правда :(
% \usepackage{pscyr}

% ----------------------------------------
% Стилевые настройки

\usepackage{boldline} % жирная линия после таблиц (чтобы не было ошибок, этот пакет должен подключаться именно тут!)
\floatsetup[table]{style=Plaintop,floatrowsep=qquad} % настройка оформления таблиц
\setlist[enumerate,itemize]{leftmargin=5mm,itemindent=10mm,itemsep=0mm,
listparindent=0em,labelsep=2mm,topsep=2mm,labelwidth=4mm} % настройки списков

\setlength{\columnsep}{0.5cm} % расстояние между колонками
\setlength{\parskip}{1pt} % расстояние до текста от колонтитула

%\usepackage{titlesec} % управление оформлением section
%\renewcommand{\thesection}{\Roman{section}}
%\titleformat{\section}[block]{\bfseries\large}{\thesection.}{5pt}{}

% Настройка совместимости
\pgfplotsset{compat=1.18}

% ----------------------------------------
% Настройки полей
\geometry{
	left=10mm,
	top=10mm,
	right=10mm,
	bottom=15mm,
	marginparsep=0mm,
	marginparwidth=0mm,
	headheight=0pt,
	headsep=0pt,
	footskip=20pt}

% ----------------------------------------
% Настройки колонтитулов и нумерации страниц
\pagenumbering{arabic}



\newcounter{ntask}
\setcounter{ntask}{0}


\newcommand{\arsh}{\mathrm{arsh} \,\,}
\newcommand{\arch}{\mathrm{arch} \,\,}
\newcommand{\arth}{\mathrm{arth} \,\,}
\newcommand{\arcth}{\mathrm{arcth} \,\,}
\renewcommand{\Re}{\operatorname{Re} \,}
\newcommand{\EDS}{\mathscr{E}}
\newcommand{\diffract}[1]{\frac{\mathrm{d}#1}{\mathrm{d}t}}

\newcommand{\kHz}{~\mathrm{kHz}}
\newcommand{\GHz}{~\mathrm{GHZ}}
\newcommand{\us}{~\mathrm{\mu s}}
\addto\captionsrussian{\def\refname{Источники}}

\begin{document}
\title{\begin{center}Лабораторная работа № 5.1.2\end{center}
Исследование эффекта Комптона}
\author{Струков О. И. \\ Б04-404}
\date{}

\thispagestyle{empty}
\pagenumbering{gobble}
\maketitle
\newpage
\pagenumbering{arabic}




\section*{Теоретическая часть}

\subsection*{1. Корпускулярно-волновой дуализм и эффект Комптона}
Рассеяние электромагнитного излучения в веществе наглядно демонстрирует двойственную природу света. Если для длинноволнового излучения успешно применяется классическая волновая теория, то при рассеянии рентгеновских и $\gamma$-лучей она не может объяснить появление в спектре рассеянного излучения дополнительной длинноволновой компоненты, которой не было в первичном пучке. 

Это явление, открытое А. Комптоном, получило название \textbf{эффекта Комптона} --- увеличения длины волны рассеянного излучения по сравнению с падающим. Эффект легко интерпретируется, если рассматривать $\gamma$-излучение как поток квантов (фотонов) с энергией $E = \hbar\omega$ и импульсом $p = \hbar\omega/c$. Эффект Комптона представляет собой результат упругого соударения фотона со свободным электроном.

\subsection*{2. Кинематика комптоновского рассеяния}
Рассмотрим столкновение $\gamma$-кванта с электроном, который до соударения покоился. Энергия покоя электрона равна $mc^2$. Пусть налетающий квант имел энергию $\hbar\omega_0$ и импульс $\hbar\omega_0/c$. После столкновения электрон приобретает энергию $\gamma mc^2$ и импульс $\gamma mv$, а $\gamma$-квант рассеивается на угол $\theta$, его энергия и импульс становятся равными $\hbar\omega_1$ и $\hbar\omega_1/c$ соответственно.

% === МЕСТО ДЛЯ ВСТАВКИ СХЕМЫ ===
\begin{figure}[htbp]
    \centering
    \includegraphics[width=0.5\textwidth]{1.jpg}
    % \fbox{\parbox{0.6\textwidth}{\centering\vspace{2cm} \textit{Вставьте сюда Рис. 1 из методички: Векторная диаграмма рассеяния $\gamma$-кванта на электроне} \vspace{2cm}}}
    \caption{Векторная диаграмма рассеяния $\gamma$-кванта на электроне.}
    \label{fig:vector}
\end{figure}

Запишем законы сохранения энергии и импульса для этого процесса:
\begin{equation}
    mc^2 + \hbar\omega_0 = \gamma mc^2 + \hbar\omega_1,
\end{equation}
\begin{equation}
    \frac{\hbar\omega_0}{c} = \gamma mv \cos\phi + \frac{\hbar\omega_1}{c} \cos\theta,
\end{equation}
\begin{equation}
    \gamma mv \sin\phi = \frac{\hbar\omega_1}{c} \sin\theta.
\end{equation}
Решая эту систему уравнений и переходя от частот к длинам волн ($\lambda = 2\pi c / \omega$), получаем выражение для изменения длины волны рассеянного излучения:
\begin{equation}
    \Delta\lambda = \lambda_1 - \lambda_0 = \frac{h}{mc}(1 - \cos\theta) = \Lambda_к(1 - \cos\theta),
\end{equation}
где $\Lambda_к = \frac{h}{mc} \approx 2,42 \cdot 10^{-10}$ см --- \textbf{комптоновская длина волны электрона}. 

Из формулы следует, что комптоновское смещение не зависит ни от длины волны первичного излучения, ни от рода вещества. В данном выводе электрон считается свободным, что справедливо для $\gamma$-квантов с энергией в сотни кэВ, так как энергия связи электронов в лёгких атомах (например, в углероде) составляет всего несколько эВ.

\subsection*{3. Взаимодействие $\gamma$-квантов с веществом}
При прохождении через вещество $\gamma$-кванты могут испытывать три основных типа взаимодействий:
\begin{enumerate}
    \item \textbf{Фотоэффект} --- квант целиком поглощается атомом, выбивая электрон с внутренней оболочки.
    \item \textbf{Комптоновское рассеяние} --- упругое столкновение кванта со слабосвязанным электроном.
    \item \textbf{Рождение электрон-позитронных пар} --- пороговый процесс, возможный только при энергии кванта $E > 2mc^2 = 1,02$ МэВ.
\end{enumerate}

В нашей работе используется источник $^{137}$Cs с энергией $\gamma$-квантов $E_0 = 662$ кэВ. Поскольку эта энергия меньше порога рождения пар, третий механизм исключён. Кроме того, при таких высоких энергиях сечение рэлеевского (несмещённого) рассеяния на атомах пренебрежимо мало по сравнению с комптоновским. Поэтому в спектре рассеянного излучения мы будем наблюдать только смещённую компоненту.

Для регистрации $\gamma$-квантов используется сцинтилляционный счётчик, состоящий из кристалла NaI(Tl) и фотоэлектронного умножителя (ФЭУ). В кристалле $\gamma$-квант рождает быстрые электроны (за счёт фотоэффекта или комптоновского рассеяния внутри самого кристалла), которые вызывают световую вспышку. Амплитуда электрического импульса на выходе ФЭУ пропорциональна энергии, поглощённой в кристалле.

% === МЕСТО ДЛЯ ВСТАВКИ СХЕМЫ ===
\begin{figure}[htbp]
    \centering
    \includegraphics[width=0.6\textwidth]{2.jpg}
    % \fbox{\parbox{0.7\textwidth}{\centering\vspace{2cm} \textit{Вставьте сюда Рис. 2 из методички: Амплитудное распределение импульсов в сцинтилляторе NaI(Tl)} \vspace{2cm}}}
    \caption{Типичный спектр импульсов от монохроматических $\gamma$-квантов в сцинтилляторе NaI(Tl).}
    \label{fig:spectrum}
\end{figure}

На рис. \ref{fig:spectrum} показано амплитудное распределение импульсов. Сплошной фон обязан комптоновскому рассеянию внутри кристалла, а чёткий пик (называемый \textbf{фотопиком} или пиком полного поглощения) возникает, когда $\gamma$-квант передаёт кристаллу всю свою энергию посредством фотоэффекта. Положение вершины фотопика однозначно связано с энергией регистрируемого излучения.

\subsection*{4. Метод обработки экспериментальных данных}
Основная цель работы --- экспериментальная проверка формулы Комптона. Для удобства обработки перейдём в формуле (4) от длин волн к энергиям $\gamma$-квантов. Разделив обе части на $\Lambda_к = h/mc$ и учитывая, что $E = hc/\lambda$, получим:
\begin{equation}
    \frac{1}{\varepsilon(\theta)} - \frac{1}{\varepsilon_0} = 1 - \cos\theta,
\end{equation}
где $\varepsilon_0 = E_0/mc^2$ и $\varepsilon(\theta) = E(\theta)/mc^2$ --- энергии падающего и рассеянного квантов, выраженные в единицах энергии покоя электрона $mc^2$.

В эксперименте энергия кванта измеряется не напрямую, а через номер канала $N(\theta)$ амплитудного анализатора, которому соответствует вершина фотопика. Так как номер канала пропорционален энергии ($N \propto E$), мы можем записать $\varepsilon(\theta) = A \cdot N(\theta)$, где $A$ --- неизвестный коэффициент калибровки. Подставив это в уравнение (5), получаем линейную зависимость:
\begin{equation}
    \frac{1}{N(\theta)} - \frac{1}{N(0)} = A(1 - \cos\theta).
\end{equation}
Если построить график зависимости $1/N(\theta)$ от $(1 - \cos\theta)$, экспериментальные точки должны лечь на одну прямую линию. 

Зная положение фотопика для угла $0^\circ$ (прямой пучок, $N(0)$) и для угла $90^\circ$ ($N(90)$), можно найти энергию покоя частицы, на которой происходит рассеяние, не зная точного коэффициента калибровки $A$. При $\theta = 90^\circ$ формула (5) принимает вид:
\begin{equation}
    mc^2 = E_0 \frac{N(90)}{N(0) - N(90)}.
\end{equation}
Для повышения точности в эту формулу рекомендуется подставлять значения $N(0)$ и $N(90)$, найденные из графика (пересечение прямой с осями), а не взятые непосредственно из сырых измерений. Это позволяет усреднить случайные погрешности и компенсировать возможные дрейфы коэффициента усиления ФЭУ.



\section*{Ход работы}


В ходе эксперимента измерялась зависимость номера канала фотопика $N(\theta)$ от угла рассеяния $\theta$ для $\gamma$-квантов источника $^{137}$Cs ($E_0 = 662$ кэВ). Рабочее напряжение на фотоэлектронном умножителе составляло $U = 1471$ В. Погрешность установки угла $\sigma\theta = 1^\circ$, погрешность определения номера канала $\sigma N = 5$.

Результаты измерений и промежуточные расчёты для построения линеаризованного графика представлены в таблице 1. Погрешность величины $1/N$ рассчитывалась по формуле $\sigma(1/N) = \sigma N / N^2$.

\begin{table}[H]
\centering
\label{tab:compton_data}
\renewcommand{\arraystretch}{1.2}
\begin{tabular}{|c|c|c|c|c|c|}
\hline
$\theta$, $^\circ$ & $\cos\theta$ & $1 - \cos\theta$ & $N(\theta)$ & $1/N(\theta)$ & $\sigma(1/N)$ \\
\hline
0   & 1,000 & 0,000 & 927 & 0,001079 & 0,000006 \\
10  & 0,985 & 0,015 & 888 & 0,001126 & 0,000006 \\
20  & 0,940 & 0,060 & 800 & 0,001250 & 0,000008 \\
30  & 0,866 & 0,134 & 745 & 0,001342 & 0,000009 \\
40  & 0,766 & 0,234 & 688 & 0,001453 & 0,000011 \\
50  & 0,643 & 0,357 & 605 & 0,001653 & 0,000014 \\
60  & 0,500 & 0,500 & 546 & 0,001832 & 0,000017 \\
70  & 0,342 & 0,658 & 490 & 0,002041 & 0,000021 \\
80  & 0,174 & 0,826 & 434 & 0,002304 & 0,000027 \\
90  & 0,000 & 1,000 & 400 & 0,002500 & 0,000031 \\
100 & -0,174& 1,174 & 356 & 0,002809 & 0,000039 \\
110 & -0,342& 1,342 & 329 & 0,003040 & 0,000046 \\
\hline
\end{tabular}
\caption{Зависимость положения фотопика от угла рассеяния}
\end{table}

По данным таблицы был построен график зависимости $1/N(\theta)$ от $(1 - \cos\theta)$:

\begin{figure}[H]
    \centering
    \includegraphics[width=0.8\textwidth]{3.pdf}
    % \fbox{\parbox{0.8\textwidth}{\centering\vspace{3cm} \textit{Вставьте сюда график: по оси X --- $1 - \cos\theta$, по оси Y --- $1/N(\theta)$ с error bars. Через точки проведена прямая.} \vspace{3cm}}}
    \caption{Зависимость $1/N(\theta)$ от $1 - \cos\theta$ и её линейная аппроксимация.}
    \label{fig:compton_graph}
\end{figure}

Методом наименьших квадратов были определены значения пересечения этой прямой с осями, что позволило найти уточнённые значения $N_{\text{наил}}$ для углов $0^\circ$ и $90^\circ$:
\begin{itemize}
    \item При $1 - \cos\theta = 0$ (ось ординат): $1/N_{\text{наил}}(0) \approx 0,00106$, откуда $N_{\text{наил}}(0) \approx 940 \pm 15$.
    \item При $1 - \cos\theta = 1$ (угол $90^\circ$): $1/N_{\text{наил}}(90) \approx 0,00260$, откуда $N_{\text{наил}}(90) \approx 385 \pm 15$.
\end{itemize}

Используя эти значения, рассчитаем энергию покоя частицы, на которой происходит рассеяние, по формуле получаем:
\begin{equation}
    mc^2 = E_0 \frac{N_{\text{наил}}(90)}{N_{\text{наил}}(0) - N_{\text{наил}}(90)} \approx 459 \text{ кэВ}
\end{equation}

% Подставляем числа ($E_0 = 662$ кэВ):
% \begin{equation}
%     mc^2 = 662 \cdot \frac{385}{940 - 385} = 662 \cdot \frac{385}{555} \approx 459 \text{ кэВ}
% \end{equation}

Оценим относительную погрешность результата:%. Основную ошибку вносит графическое определение $N_{\text{наил}}$:
\begin{equation}
    \frac{\sigma(mc^2)}{mc^2} \approx \sqrt{ \left(\frac{\sigma N_{90}}{N_{90}}\right)^2 + \left(\frac{\sqrt{\sigma N_0^2 + \sigma N_{90}^2}}{N_0 - N_{90}}\right)^2 } \approx 0,055.
\end{equation}
Абсолютная погрешность: $\sigma(mc^2) = 459 \cdot 0,055 \approx 25$ кэВ.

Итоговый результат:
\begin{equation}
    mc^2 = 460 \pm 25 \text{ кэВ}.
\end{equation}

\subsection*{Вывод}
Табличное значение энергии покоя электрона составляет $511$ кэВ. Полученный нами результат ($460 \pm 25$ кэВ) немного занижен, но табличное значение попадает в диапазон $2\sigma$. Небольшое расхождение может быть вызвано неточностью калибровки спектрометра, краевыми эффектами в кристалле NaI(Tl) или погрешностью геометрии рассеяния. Тем не менее, линейный характер зависимости $1/N(\theta)$ от $1 - \cos\theta$ наглядно подтверждает справедливость формулы Комптона.



















\end{document}